\section{Erd\H{o}s Problem 164: the weighted sum of a primitive set}

\subsection{1) Statement and scope}
A set $A\subseteq\{2,3,\ldots\}$ is primitive if distinct members never
divide one another. The assertion is
\[
 \sum_{a\in A}\frac1{a\log a}
 \leq \sum_{p\ {\rm prime}}\frac1{p\log p}.
\]
It was proved by Lichtman. The following reconstruction uses the von
Mangoldt-chain approach of Alexeev and collaborators (2026), with one
analytic inequality stated explicitly as an external input. The
probability argument and the exceptional prime-power adjustment are
proved here. Logarithms are natural.

\subsection{2) The analytic input and a downward chain}
Write $w(n)=1/(n\log n)$ and let $\Lambda$ be the von Mangoldt function.
Lemma 3.3(ii) of the cited 2026 paper states, for every real $m\geq2$,
with $x=\log m/\log2$,
\[
 \log m\sum_{q\geq2}\frac{\Lambda(q)}{q\log^2(mq)}
 \leq\frac{x}{x+1/2}.                                      \tag{1}
\]
This analytic estimate is not reproved in this note.

For a composite integer $n$ which is not a prime power, set
$T(n,n/q)=\Lambda(q)/\log n$ for each divisor $q\geq2$ of $n$.
Terms with zero numerator are ignored. The destination is always at
least $2$, and $\sum_{q\mid n}\Lambda(q)=\log n$ shows that these
probabilities sum to one. For $n=p^k$, $k\geq2$, instead put
\[
 T(p^k,p)=2/k,\qquad
 T(p^k,p^j)=1/k\quad(2\leq j\leq k-1).
\]
The empty second range when $k=2$ is allowed. These probabilities also
sum to one. Primes have no downward transitions.

Let $W(m)=\sum_{n>m}w(n)T(n,m)$. If $m$ is composite, direct substitution
gives
\[
 \frac{W(m)}{w(m)}=
 \log m\sum_{q\geq2}\frac{\Lambda(q)}{q\log^2(mq)}\leq1.
\]
If $m=p$ is prime, the doubled last transition contributes the additional
term $\sum_{k\geq2}1/(k^2p^{k-1})$. Since $x=\log_2p\leq p-1$,
\[
 \sum_{k\geq2}\frac1{k^2p^{k-1}}
 \leq\frac1{4(p-1)}\leq\frac1{2p-1}
 \leq\frac1{2x+1}.
\]
Together with (1), this yields $W(p)/w(p)\leq
x/(x+1/2)+1/(2x+1)=1$. Thus $W(m)\leq w(m)$ for every $m\geq2$.

\subsection{3) Reverse the chain and count visits}
The prime mass $P=\sum_p w(p)$ is finite. For completeness, primes in
$(2^j,2^{j+1}]$ divide $\binom{2^{j+1}}{2^j}\leq4^{2^j}$, so there
are at most $2^{j+1}/j$ of them for $j\geq1$. Their total weight is at
most $2/(j^2\log2)$; adding the prime $2$ proves convergence.

Start a random chain at a prime $p$, with probability $w(p)/P$.
From $m$ move to a proper multiple $n>m$ with probability
\[
 Q(m,n)=\frac{w(n)T(n,m)}{w(m)}.
\]
The preceding bound says that the row sum of $Q$ is at most one.
With the remaining probability the chain terminates. Thus this defines
a countable-state process; its integer states strictly increase and
each divides the next.

Let $h(n)$ be $P$ times the probability that the process ever visits
$n$. A prime can only be visited initially, so $h(p)=w(p)$. For composite
$n$, its possible predecessors are its proper divisors at least $2$.
Induction on $n$ and the fact that a state cannot be visited twice give
\[
 h(n)=\sum_{\substack{m\mid n\\2\leq m<n}}h(m)Q(m,n)
 =w(n)\sum_m T(n,m)=w(n).
\]
Every trajectory visits at most one member of a primitive set. Therefore,
by the nonnegative sum rule for expectations,
\[
 \sum_{a\in A}w(a)
 =P\,\mathbb E\bigl[\#\{\hbox{visited states in }A\}\bigr]\leq P.
\]
This includes infinite primitive sets, without an unproved exchange of
signed series. The primes themselves are primitive and attain $P$.

\subsection{4) Outcome and sources}
\textbf{Outcome: SOLVED using the stated analytic lemma.} All steps from
(1) to the desired inequality are proved above. The original resolution
and the later chain method belong to the cited authors; this is an
exposition, not a novelty or formal-verification claim.
\begin{itemize}
\item B. Alexeev, K. Barreto, Y. Li, J. D. Lichtman, L. Price,
J. I. Shah, Q. Tang, and T. Tao, \emph{Primitive sets and von Mangoldt
chains: Erd\H{o}s Problem \#1196 and beyond}, Lemma 3.3(ii) and Section 5,
\texttt{https://arxiv.org/abs/2605.00301}.
\item T. F. Bloom, Problem 164,
\texttt{https://www.erdosproblems.com/164}, checked 23 September 2026.
\end{itemize}
The percentage includes the explicitly cited analytic input.
\subsection{5) COMPLETION ESTIMATE}
\textbf{COMPLETION: 100\%}
